\section{Introduction}

Run the same six eviction policies on two real public-records API workloads
and the policies rank almost identically --- with one exception. On our
20{,}000-request Congress.gov trace, SIEVE and W-TinyLFU finish in a
photo-finish at high skew, gaps smaller than a single seed's standard
deviation. On our 20{,}000-request CourtListener trace, W-TinyLFU pulls
ahead of SIEVE by a consistent $4$--$5$\,pp across every $\alpha$ we tested.
Figure~\ref{fig:opening} (page~1) shows the split: same policies, same
$\alpha$ sweep, same five-seed aggregation with $\pm 1\sigma$ bands --- and
yet the Congress curves all collapse into a narrow band near
$\alpha = 1.2$ while the Court curves fan out into clean strata.

This paper explains why. The short answer is that it is not about the
policies at all; it is about the workloads. The raw CourtListener trace
has $\alpha = 1.03$ under MLE, a genuinely heavy-tailed access pattern.
The raw Congress.gov trace has $\alpha = 0.23$, effectively uniform. We
overlay a replay-Zipf popularity distribution on both traces to make
policy comparison meaningful at all, but the raw $\alpha$ gap still bleeds
through --- primarily because Court carries a single $462$\,KB outlier
object (a $1{,}984\times$ multiple of the median document size) that
admission-filter-based policies like W-TinyLFU handle gracefully and
recency-only policies like SIEVE do not. Congress has no such outlier,
so the admission-filter advantage evaporates and SIEVE catches up at
high $\alpha$. Across the full Welch's-t-tested grid, W-TinyLFU dominates
the miss-ratio dimension in 11 of 12 cells (six cache fractions $\times$
two workloads); SIEVE is the only other winner and only at one Congress
cell. The mechanism deep-dive in \S\ref{sec:mechanism} unpacks this
exact story, grounded in the raw-$\alpha$ and size-distribution numbers
behind Figure~\ref{fig:opening}.

We frame methodology choices in first person (``We chose replay-Zipf
because the raw trace's near-zero temporal locality makes pure-replay
comparisons meaningless'') and results in third person (``W-TinyLFU
wins 11 of 12 cells''), matching the standard CS-class-report register.

\paragraph{Contributions.} This paper makes four contributions:
\begin{enumerate}[leftmargin=*,itemsep=0pt,topsep=2pt]
  \item \textbf{Two real legislative workloads characterized side-by-side}
    (\S\ref{sec:workload}): Congress.gov bill metadata and CourtListener
    court documents, each $\approx 20$\,K requests, with raw $\alpha$,
    OHW ratio, size distribution, and unique-key fraction reported
    jointly so readers can see where the workloads agree and diverge.
  \item \textbf{Five-seed Welch's-t-backed policy comparison with CI
    bands} (\S\ref{sec:policies}): every claim of the form ``policy A
    beats policy B on workload W at cache fraction $f$'' is backed by
    $p < 0.05$ across 5 seeds, with $\pm 1\sigma$ bands plotted so the
    reader sees the spread, not just the mean.
  \item \textbf{A SIEVE-vs-W-TinyLFU mechanism explanation grounded in
    raw $\alpha \times$ size distribution} (\S\ref{sec:mechanism}): we
    show that the SIEVE/W-TinyLFU tie on Congress and the $4$--$5$\,pp
    gap on Court both follow from two workload traits rather than any
    property of the policies themselves.
  \item \textbf{A practitioner decision tree}
    (\S\ref{sec:conclusion}): given raw $\alpha$ and the presence or
    absence of catastrophic size outliers in a public-records API
    workload, which of the six policies to reach for first.
\end{enumerate}

The rest of the paper is organized as follows. \S\ref{sec:related}
situates the work against SIEVE (\cite{zhang2024sieve}), S3-FIFO
(\cite{yang2023fifo}), TinyLFU (\cite{einziger2017tinylfu}), Caffeine
(\cite{caffeine2024}), and SHARDS (\cite{waldspurger2015shards}).
\S\ref{sec:workload} characterizes both traces. \S\ref{sec:policies}
presents the six-policy miss-ratio comparison.
\S\ref{sec:mechanism} is the paper's analytical centerpiece.
\S\ref{sec:shards} validates our SHARDS MRC construction at four
sampling rates. \S\ref{sec:ablations} isolates three mechanism claims
through controlled ablations. \S\ref{sec:byte-mrc} flags a
byte-miss-ratio variance warning readers should heed before trusting
per-seed byte numbers. \S\ref{sec:conclusion} closes with the decision
tree, limitations, and v2 directions.
